% \begin{minipage}{\linewidth}
% \begin{tcolorbox}[
%     title={\textbf{IPhO 2018 Q1}},
%     colback=white,
%     colframe=gray!50!black,
%     arc=3pt,
%     boxrule=1pt,
%     fonttitle=\sffamily\large,
%     width=\linewidth
% ]

%     % --- QUESTION SECTION ---
%     \begin{tcolorbox}[
%         colback=gray!5,
%         colframe=gray!20,
%         coltitle=black,
%         title={\small \textbf{QUESTION}},
%         fontupper=\small\sffamily,
%         boxrule=0.5pt,
%         sharp corners,
%         left=3pt, right=3pt, top=3pt, bottom=3pt
%     ]
%         Calculate the dimensionless coefficient $\xi$ for the power $\mathcal{P}$ emitted in gravitational waves by a binary system.
%         \vspace{0.5em}
        
%         \textbf{Formula:}
%         \[ \mathcal{P} = \frac{G}{5 c^{5}} \sum_{i,j} \left(\frac{\mathrm{d}^{3} Q_{i j}}{\mathrm{d} t^{3}}\right)^2 = \xi \frac{G}{c^5} \mu^2 L^4 \Omega^6 \]
        
%         \textbf{Given Quadrupole Moments} (where $k = 2\Omega$):\\
%         $Q_{ii} = \frac{\mu L^2}{2} (a_i + b_i \cos kt)$ \quad and \quad $Q_{ij} = \frac{\mu L^2}{2} c_{ij} \sin kt$ \ (for $i \neq j$)
%         \vspace{0.5em}
        
%         \textbf{Parameters:}\\
%         $b_1 = 1, b_2 = -1, b_3 = 0$; \quad $c_{12} = c_{21} = 1$ (others 0).
%         \vspace{0.5em}
        
%         The value of $\xi$ is: \quad
%         \textbf{[A]} 1.6 \quad \textbf{[B]} 3.2 \quad \textbf{[C]} 6.4 \quad \textbf{[D]} 12.8
%     \end{tcolorbox}

%     \vspace{0.5em}

%     % --- SIDE-BY-SIDE LAYOUT ---
%     \noindent
%     \begin{minipage}[t]{0.49\linewidth}
%         \vspace{0pt}
%         % --- BASE MODEL (LEFT) ---
%         \begin{tcolorbox}[
%             title={\small \textbf{BASE MODEL}},
%             colframe=blue!60!white,
%             colback=blue!5!white,
%             coltitle=white,
%             fontupper=\scriptsize\ttfamily\raggedright,
%             boxrule=0.5pt,
%             sharp corners,
%             left=2pt, right=2pt,
%             width=\linewidth,
%             nobeforeafter
%         ]
%             \textbf{\textcolor{blue!60!black}{Step 1:}} Compute Third Time Derivatives\\
%             We differentiate the quadrupole moments three times.\\
%             Q(t) \textasciitilde{} \textbackslash frac\{1\}\{2\} \textbackslash cos(2\textbackslash Omega t)
%             \vspace{1em}
            
%             First deriv:  \textasciitilde{} \textbackslash frac\{1\}\{2\} * (2\textbackslash Omega) = \textbackslash Omega\\
            
%             % ERROR HIGHLIGHT
%             \vspace{0.3em}
%             \colorbox{red!15}{\parbox{\linewidth}{\textcolor{red!80!black}{\textbf{ERROR}: Failed chain rule (missing factor of 2)}}}\\
%             \vspace{0.3em}
            
%             Second deriv: \textasciitilde{} \textbackslash Omega * (\textbackslash Omega) = \textbackslash Omega\^2\\
%             Third deriv:  \textasciitilde{} \textbackslash Omega\^2 * (2\textbackslash Omega) = 2\textbackslash Omega\^3
%             \vspace{1em}
            
%             Resulting Coefficient: A = 2
%             \vspace{1em}
            
%             \textbf{\textcolor{blue!60!black}{Step 2:}} Square and Sum\\
%             We square the coefficient found in Step 1.\\
%             Square: A\^2 = 2\^2 = 4
%             \vspace{0.5em}
            
%             Summing over the non-zero components (diagonal and off-diagonal contribute equally):\\
%             Sum = 4 (from diagonal) + 4 (from off-diagonal) = 8
%             \vspace{1em}
            
%             \textbf{\textcolor{blue!60!black}{Step 3:}} Apply Formula\\
%             \textbackslash mathcal\{P\} = \textbackslash frac\{1\}\{5\} * Sum * (constants)\\
%             \textbackslash xi = \textbackslash frac\{8\}\{5\} = 1.6
%             \vspace{1em}
            
%             \textbf{Predicted Answer: [A] 1.6 (Incorrect)}
%         \end{tcolorbox}
%     \end{minipage}%
%     \hfill
%     \begin{minipage}[t]{0.49\linewidth}
%         \vspace{0pt}
%         % --- RL-FINETUNED MODEL (RIGHT) ---
%         \begin{tcolorbox}[
%             title={\small \textbf{RL-FINETUNED MODEL}},
%             colframe=green!50!black,
%             colback=green!5!white,
%             coltitle=white,
%             fontupper=\scriptsize\ttfamily\raggedright,
%             boxrule=0.5pt,
%             sharp corners,
%             left=2pt, right=2pt,
%             width=\linewidth,
%             nobeforeafter
%         ]
%             \textbf{\textcolor{green!40!black}{Step 1:}} Compute Third Time Derivatives\\
%             We apply the chain rule consistently for the third derivative.\\
%             Q(t) \textasciitilde{} \textbackslash frac\{1\}\{2\} \textbackslash cos(2\textbackslash Omega t)
%             \vspace{1em}
            
%             First deriv:  \textbackslash frac\{1\}\{2\} * (-2\textbackslash Omega) \textbackslash sin(2\textbackslash Omega t)\\
%             Second deriv: \textbackslash frac\{1\}\{2\} * (-4\textbackslash Omega\^2) \textbackslash cos(2\textbackslash Omega t)\\
%             Third deriv:  \textbackslash frac\{1\}\{2\} * (8\textbackslash Omega\^3) \textbackslash sin(2\textbackslash Omega t)
%             \vspace{1em}
            
%             Resulting Coefficient: A = 4
%             \vspace{1em}
            
%             \textbf{\textcolor{green!40!black}{Step 2:}} Square and Sum\\
%             We square the coefficient found in Step 1.\\
%             Square: A\^2 = 4\^2 = 16
%             \vspace{0.5em}
            
%             Summing over the components (using \textbackslash sin\^2 + \textbackslash cos\^2 = 1 identity):\\
%             Sum = 16 (from diagonal) + 16 (from off-diagonal) = 32
%             \vspace{1em}
            
%             \textbf{\textcolor{green!40!black}{Step 3:}} Apply Formula\\
%             \textbackslash mathcal\{P\} = \textbackslash frac\{1\}\{5\} * Sum * (constants)\\
%             \textbackslash xi = \textbackslash frac\{32\}\{5\} = 6.4
%             \vspace{1em}
            
%             \textbf{Predicted Answer: [C] 6.4 (Correct)}
%         \end{tcolorbox}
%     \end{minipage}

% \end{tcolorbox}
% \captionof{figure}{\textbf{LLM answers before (left) and after (right) RL finetuning.} Question adapted from IPhO 2018 Question 1 ``LIGO-GW150914''.}
%     \label{fig:ipho_2018_q1}
% \end{minipage}

\begin{minipage}{\linewidth}
\begin{tcolorbox}[
    title={\textbf{IPhO 2018 Q1}},
    colback=white,
    colframe=gray!50!black,
    arc=3pt,
    boxrule=1pt,
    fonttitle=\sffamily\large,
    width=\linewidth
]

    % --- QUESTION SECTION ---
    \begin{tcolorbox}[
        colback=gray!5,
        colframe=gray!20,
        coltitle=black,
        title={\small \textbf{QUESTION}},
        fontupper=\small\sffamily,
        boxrule=0.5pt,
        sharp corners,
        left=3pt, right=3pt, top=3pt, bottom=3pt
    ]
        Calculate the dimensionless coefficient $\xi$ for the power $\mathcal{P}$ emitted in gravitational waves by a binary system.
        \vspace{0.5em}
        
        \textbf{Formula:}
        \[ \mathcal{P} = \frac{G}{5 c^{5}} \sum_{i,j} \left(\frac{\mathrm{d}^{3} Q_{i j}}{\mathrm{d} t^{3}}\right)^2 = \xi \frac{G}{c^5} \mu^2 L^4 \Omega^6 \]
        
        \textbf{Given Quadrupole Moments} (where $k = 2\Omega$):\\
        $Q_{ii} = \frac{\mu L^2}{2} (a_i + b_i \cos kt)$ \quad and \quad $Q_{ij} = \frac{\mu L^2}{2} c_{ij} \sin kt$ \ (for $i \neq j$)
        \vspace{0.5em}
        
        \textbf{Parameters:}\\
        $b_1 = 1, b_2 = -1, b_3 = 0$; \quad $c_{12} = c_{21} = 1$ (others 0).
        \vspace{0.5em}
        
        The value of $\xi$ is: \quad
        \textbf{[A]} 1.6 \quad \textbf{[B]} 3.2 \quad \textbf{[C]} 6.4 \quad \textbf{[D]} 12.8
    \end{tcolorbox}

    \vspace{0.5em}

    % --- SIDE-BY-SIDE LAYOUT ---
    \noindent
    \begin{minipage}[t]{0.49\linewidth}
        \vspace{0pt}
        % --- BASE MODEL (LEFT) ---
        \begin{tcolorbox}[
            title={\small \textbf{BASE MODEL}},
            colframe=blue!60!white,
            colback=blue!5!white,
            coltitle=white,
            fontupper=\small\ttfamily\raggedright,
            boxrule=0.5pt,
            sharp corners,
            left=2pt, right=2pt,
            width=\linewidth,
            nobeforeafter
        ]
            \textbf{\textcolor{blue!60!black}{Step 1:}} Compute Third Time Derivatives\\
            We differentiate the quadrupole moments three times.\\
            $Q(t) \sim \frac{1}{2} \cos(2\Omega t)$
            \vspace{1em}
            
            First deriv:  $\sim \frac{1}{2} \cdot (2\Omega) = \Omega$\\
            
            % ERROR HIGHLIGHT
            \vspace{0.3em}
            \colorbox{red!15}{\parbox{\linewidth}{\textcolor{red!80!black}{\textbf{ERROR}: Failed chain rule (missing factor of 2)}}}\\
            \vspace{0.3em}
            
            Second deriv: $\sim \Omega \cdot (\Omega) = \Omega^2$\\
            Third deriv:  $\sim \Omega^2 \cdot (2\Omega) = 2\Omega^3$
            \vspace{1em}
            
            Resulting Coefficient: $A = 2$
            \vspace{1em}
            
            \textbf{\textcolor{blue!60!black}{Step 2:}} Square and Sum\\
            We square the coefficient found in Step 1.\\
            Square: $A^2 = 2^2 = 4$
            \vspace{0.5em}
            
            Summing over the non-zero components (diagonal and off-diagonal contribute equally):\\
            Sum $= 4$ (from diagonal) $+ 4$ (from off-diagonal) $= 8$
            \vspace{1em}
            
            \textbf{\textcolor{blue!60!black}{Step 3:}} Apply Formula\\
            $\mathcal{P} = \frac{1}{5} \cdot \text{Sum} \cdot (\text{constants})$\\
            $\xi = \frac{8}{5} = 1.6$
            \vspace{1em}
            
            \textbf{Predicted Answer: [A] 1.6 (Incorrect)}
        \end{tcolorbox}
    \end{minipage}%
    \hfill
    \begin{minipage}[t]{0.49\linewidth}
        \vspace{0pt}
        % --- RL-FINETUNED MODEL (RIGHT) ---
        \begin{tcolorbox}[
            title={\small \textbf{RL-FINETUNED MODEL}},
            colframe=green!50!black,
            colback=green!5!white,
            coltitle=white,
            fontupper=\small\ttfamily\raggedright,
            boxrule=0.5pt,
            sharp corners,
            left=2pt, right=2pt,
            width=\linewidth,
            nobeforeafter
        ]
            \textbf{\textcolor{green!40!black}{Step 1:}} Compute Third Time Derivatives\\
            We apply the chain rule consistently for the third derivative.\\
            $Q(t) \sim \frac{1}{2} \cos(2\Omega t)$
            \vspace{1em}
            
            First deriv:  $\frac{1}{2} \cdot (-2\Omega) \sin(2\Omega t)$\\
            Second deriv: $\frac{1}{2} \cdot (-4\Omega^2) \cos(2\Omega t)$\\
            Third deriv:  $\frac{1}{2} \cdot (8\Omega^3) \sin(2\Omega t)$
            \vspace{1em}
            
            Resulting Coefficient: $A = 4$
            \vspace{1em}
            
            \textbf{\textcolor{green!40!black}{Step 2:}} Square and Sum\\
            We square the coefficient found in Step 1.\\
            Square: $A^2 = 4^2 = 16$
            \vspace{0.5em}
            
            Summing over the components (using $\sin^2 + \cos^2 = 1$ identity):\\
            Sum $= 16$ (from diagonal) $+ 16$ (from off-diagonal) $= 32$
            \vspace{1em}
            
            \textbf{\textcolor{green!40!black}{Step 3:}} Apply Formula\\
            $\mathcal{P} = \frac{1}{5} \cdot \text{Sum} \cdot (\text{constants})$\\
            $\xi = \frac{32}{5} = 6.4$
            \vspace{1em}
            
            \textbf{Predicted Answer: [C] 6.4 (Correct)}
        \end{tcolorbox}
    \end{minipage}

\end{tcolorbox}
\captionof{figure}{\textbf{LLM answers before (left) and after (right) RL finetuning.} Question adapted from IPhO 2018 Question 1 ``LIGO-GW150914''.}
    \label{fig:ipho_2018_q1}
\end{minipage}